\section{Discussion}
\label{sec:implications}

\paragraph{What ``when'' means in this paper.}
We do not claim that retained-behavior damage risk is a universal property
of an example or that one scalar explains the causal mechanism of collateral
damage. Loss Susceptibility is the gradient-based first-order responsiveness
of the behavior loss within the declared block; Representation Proximity
describes
exposure to the current forget request.
The sealed score combines these two prospective coordinates for a declared
model, request, and parent regime. RQ1 determines empirically whether that
ordering anticipates later damage and where it breaks down.

\paragraph{Why prediction and protection are separate.}
A useful ranking need not identify examples that a fixed repair operator can
protect efficiently. Conversely, a repair pool may work because its examples
generalize well under rehearsal even when the original rank correlation is
modest. RQ3 therefore holds the operator, quota, compute budget, and
feasibility region fixed and changes only the selector. Improvement on the
score-selected discovery pool is insufficient; decision value requires lower
damage on the untouched audit fold and net benefit over no repair.

\paragraph{Operational role.}
The framework is most useful when exhaustive retained-behavior auditing or
repair is too expensive. It narrows an enumerated candidate universe to
behaviors worth inspecting or protecting before target outcomes are available.
This benefit is amortized only when suitable development trajectories exist to
calibrate the scalar mixture and repair contract. Invalid profiles,
gate-non-reaching parents, unsupported damage tails, and infeasible repairs
are therefore substantive deployment boundaries, not missing data to be
silently removed.
